\section{Заключение}
\label{sec:conclusion}

В данной выпускной квалификационной работе математическая задача переноса вероятностной меры $T_{\#}\pi_N = \pi_C$ для распределений нативных и контрастных КТ-изображений исследована средствами теории условного флоу матчинга. Показано, что при линейном интерполянте~\eqref{eq:linear_interp} соответствующее условное поле скоростей постоянно по времени, кривизна траектории нулевая, а кинетическая энергия минимальна по Бенаму--Бренье. Это даёт теоретическое обоснование одношагового численного интегрирования (\S\ref{sec:ode_methods}) и симметрии единственной сети для обоих направлений трансляции (\S\ref{sec:notation}). Практической реализацией служит \textbf{Medical Flow Matching~(MFM)}~--- конвейер обучения и инференса, а параметризацией поля скоростей~--- предложенная архитектура \textbf{TimeResNet} (U-Net с поблочной аффинной модуляцией временного сигнала и свёрточным самовниманием на боттлнеке).

\paragraph{Итоги работы.}

\textit{Теоретические результаты.} Из симметрии линейного интерполянта следует, что единственная сеть обеспечивает трансляцию в обоих направлениях без штрафа по качеству~--- экспериментально подтверждено разницей 0.03\,HU между двунаправленным и однонаправленным вариантами (таблица~\ref{tab:unidirectional}). Оценка ошибки метода Эйлера~\eqref{eq:euler_error} объясняет, почему одношаговый инференс достаточен: кривизна выученного поля мала при линейном интерполянте.

\textit{Качество реконструкции.} На отложенной выборке MFM достигает MAE\,=\,4.25\,HU, SSIM\,=\,0.994, PSNR\,=\,42.19\,дБ~--- лучших значений среди всех рассмотренных методов. По сравнению со SMILE MAE ниже в 7.1~раза. На клинически значимых структурах разрыв ещё значительнее: на аорте при трансляции нативное~$\to$~контрастное MFM обеспечивает MAE в~15~раз ниже, чем SMILE.

\textit{Скорость и эффективность.} Даже в честном сравнении с ускоренным сэмплером DDIM-50 MFM работает в~40~раз быстрее (0.125\,с против 5.03\,с) и в~5~раз точнее по MAE. При этом DDIM-50 деградирует по SSIM относительно полного DDPM.

\textit{Межосевая согласованность.} Несмотря на 2D-обработку, $\Delta$-SSIM MFM равен 0.0001~--- фактически нулевая межосевая рассогласованность, что подтверждается метрикой 3D-TotalSeg-LPIPS~=~0.007.

\textit{Практическая применимость.} nnUNet\,v2, обученный на MFM-переведённых нативных сканах, достигает Dice~=~0.926 против 0.966 на реальных данных~--- потеря лишь в 4 процентных пункта.
